% Block 4 of lista-reais-lectures-01-05 (input by the drivers in Listas/).
% Order and ratings: lista-reais-lectures-01-05/order.md.
\bloco{4}{Credible threats and subgame perfection}{Lecture 4}

% ---------------------------------------------------------------------------
\exercicio[1]{Repeated-game definitions}{Foundations}{Gibbons (1992), \emph{Game Theory for Applied
Economists}, Ch.~2 ``Dynamic Games of Complete Information'', \S2.6 Problems (Section 2.3), Problem
2.12, p.~135 [PDF 142]. No answer in the book. Reproduced verbatim.}

\rot{2.12.}What is a strategy in a repeated game? What is a subgame in a repeated game? What is a
subgame-perfect Nash equilibrium?
\resolespaco

% ---------------------------------------------------------------------------
\exercicio[1]{Entering an Industry}{Foundations}{Tadelis (2013), \emph{Game Theory: An Introduction},
Ch.~7 ``Preliminaries'', Exercise 7.6 ``Entering an Industry'', pp.~148--149 [PDF 167--168]. No
answer in the book. Reproduced verbatim.}

\rot{7.6}\textbf{Entering an Industry:} A firm (player 1) is considering entering an established
industry with one incumbent firm (player 2). Player 1 must choose whether or not to enter the
industry. If player 1 enters the industry then player 2 can either accommodate the entry or fight
the entry by waging a price war. Player 1's most-preferred outcome is entering with player 2 not
fighting, and its least-preferred outcome is entering with player 2 fighting. Player 2's
most-preferred outcome is player 1 not entering, and its least-preferred outcome is player 1
entering with player 2 fighting.
\begin{partes}
\item Model this as an extensive-form game tree (choose payoffs that represent the preferences).
\item How many pure strategies does each player have?
\item Find all the Nash equilibria of this game.
\end{partes}
\resolespaco

% ---------------------------------------------------------------------------
\exercicio[2]{Entry Deterrence 1}{Foundations}{Tadelis (2013), \emph{Game Theory: An Introduction},
Ch.~8 ``Credibility and Sequential Rationality'', Exercise 8.8 ``Entry Deterrence 1'',
pp.~171--172 [PDF 190--191]. No answer in the book. Reproduced verbatim.}

\rot{8.8}\textbf{Entry Deterrence 1:} NSG Corporation is considering entry into the local phone
market in the San Francisco Bay Area. The incumbent, S\&P Corporation, predicts that a price war
will result if NSG enters. If NSG stays out, S\&P earns monopoly profits valued at \$10 million (net
present value of profits) while NSG earns zero. If NSG enters, it must incur irreversible entry
costs of \$2 million. If there is a price war, each firm earns \$1 million. S\&P always has the
option of accommodating entry (i.e., not starting a price war). In such a case both firms earn \$4
million (net present value). Suppose that the timing is such that NSG first has to choose whether or
not to enter the market. Then S\&P decides whether to ``accommodate entry'' or ``engage in a price
war.'' What is the subgame-perfect equilibrium outcome to this sequential game? Set up a game tree.
\resolespaco

% ---------------------------------------------------------------------------
\exercicio[3]{Investment in the Future}{Exam level}{Tadelis (2013), \emph{Game Theory: An
Introduction}, Ch.~8 ``Credibility and Sequential Rationality'', Exercise 8.6 ``Investment in the
Future'', pp.~170--171 [PDF 189--190]. No answer in the book. Reproduced verbatim.}

\rot{8.6}\textbf{Investment in the Future:} Consider two firms that play a Cournot competition game
with demand $p = 100 - q$ and costs for each firm given by $c_i(q_i) = 10q_i$. Imagine that before
the two firms play the Cournot game firm 1 can invest in cost reduction. If it invests the costs of
firm 1 will drop to $c_1(q_1) = 5q_1$. The cost of investment is $F > 0$. Firm 2 does not have this
investment opportunity.
\begin{partes}
\item Find the value $F^*$ for which the unique subgame-perfect equilibrium involves firm 1
investing.
\item Assume that $F > F^*$. Find a Nash equilibrium of the game that is not subgame perfect.
\end{partes}
\resolespaco

% ---------------------------------------------------------------------------
\exercicio[3]{Stackelberg with Two Firms}{Exam level}{Muñoz-Garcia \& Toro-Gonzalez (2019),
\emph{Strategy and Game Theory: Practice Exercises with Answers}, 2nd ed., Ch.~4
``Sequential-Move Games with Complete Information'', Exercise 11 ``Stackelberg with Two
Firms''$^{\mathrm{A}}$, p.~180 [PDF 192]. Worked answer in the book, right after the statement.
Reproduced verbatim (statement only).}

Consider a leader and a follower in a Stackelberg game of quantity competition. Firms face an
inverse demand curve $p(Q) = 1 - Q$, where $Q = q_L + q_F$ denotes aggregate output. The leader faces
a constant marginal cost $c_L > 0$ while the follower's constant marginal cost is $c_F > 0$, where
$1 > c_F \ge c_L$, indicating that the leader has access to a more efficient technology than the
follower.

\textbf{Part (a)} Find the follower's best response function, $BRF_F$, i.e., $q_F(q_L)$.

\textbf{Part (b)} Determine each firm's output strategy in the subgame perfect equilibrium (SPNE) of
this sequential-move game.

\textbf{Part (c)} Under which conditions on $c_L$ can you guarantee that both firms produce strictly
positive output levels?

\textbf{Part (d)} Assuming cost symmetry, i.e., $c_L = c_F = c$, determine the aggregate equilibrium
output level in the Stackelberg game, and its associated equilibrium price. Then compare them with
those arising in (1) a monopoly; in (2) the Cournot game of quantity competition; and (3) in the
Bertrand game of price competition (which, under cost symmetry, is analogous to a perfectly
competitive industry) [\emph{Hint}: For a more direct and visual comparison, depict the inverse
demand curve $p(Q) = 1 - Q$, and locate each of these four aggregate output levels in the horizontal
axis (with its associated prices in the vertical axis).].
\resolespaco

% ---------------------------------------------------------------------------
\exercicio[3]{Resident Assistant Selection with a Twist}{Exam level}{Tadelis (2013), \emph{Game Theory:
An Introduction}, Ch.~8 ``Credibility and Sequential Rationality'', Exercise 8.11 ``Resident
Assistant Selection with a Twist'', pp.~172--173 [PDF 191--192]. No answer in the book. Reproduced
verbatim.}

\rot{8.11}\textbf{Resident Assistant Selection with a Twist:} Two staff managers in the
$\Pi\mathrm{B}\Phi$ sorority, the house manager (player 1) and kitchen manager (player 2), must
select a resident assistant from a pool of three candidates: $\{a, b, c\}$. Player 1 prefers $a$ to
$b$, and $b$ to $c$. Player 2 prefers $b$ to $a$, and $a$ to $c$. The process that is imposed on them
is as follows: First, the house manager vetoes one of the candidates and announces the veto to the
central office for staff selection and to the kitchen manager. Next the kitchen manager vetoes one
of the remaining two candidates and announces it to the central office. Finally the director of the
central office assigns the remaining candidate to be a resident assistant at $\Pi\mathrm{B}\Phi$.
\begin{partes}
\item Model this as an extensive-form game (using a game tree) in which a player's most-preferred
candidate gives a payoff of 2, the second most-preferred candidate gives a payoff of 1, and the last
candidate gives 0.
\item Find the subgame-perfect equilibrium for this game. Is it unique?
\item Are there Nash equilibria that are not subgame-perfect equilibria?
\item Now assume that before the two players play the game, player 2 can send an alienating e-mail
to one of the candidates, which would result in that candidate withdrawing her application. Would
player 2 choose to do this, and if so, with which candidate?
\end{partes}
\resolespaco

% ---------------------------------------------------------------------------
\exercicio[4]{Industry Leader}{Exam level}{Tadelis (2013), \emph{Game Theory: An Introduction}, Ch.~8
``Credibility and Sequential Rationality'', Exercise 8.4 ``Industry Leader'', p.~170 [PDF 189]. No
answer in the book. Reproduced verbatim. (Gibbons Problem 2.6, p.~133, poses the same game and asks
only for the subgame-perfect outcome; it is not repeated in this list.)}

\rot{8.4}\textbf{Industry Leader:} Three oligopolists operate in a market with inverse demand given
by $P(Q) = a - Q$, where $Q = q_1 + q_2 + q_3$, and $q_i$ is the quantity produced by firm $i$. Each
firm has a constant marginal cost of production, $c$, and no fixed cost. The firms choose their
quantities dynamically as follows: Firm 1, which is the industry leader, chooses $q_1 \ge 0$. Firms
2 and 3 observe $q_1$ and then simultaneously choose $q_2$ and $q_3$, respectively.
\begin{partes}
\item How many proper subgames does this dynamic game have? Explain briefly.
\item Is it a game of perfect or imperfect information? Explain briefly.
\item What is the subgame-perfect equilibrium of this game? Show that it is unique.
\item Find a Nash equilibrium that is not a subgame-perfect equilibrium.
\end{partes}
\resolespaco

% ---------------------------------------------------------------------------
\exercicio[4]{Junk Mail Advertising}{Exam level}{Tadelis (2013), \emph{Game Theory: An Introduction},
Ch.~8 ``Credibility and Sequential Rationality'', Exercise 8.13 ``Junk Mail Advertising'',
pp.~173--174 [PDF 192--193]. No answer in the book. Reproduced verbatim.}

\rot{8.13}\textbf{Junk Mail Advertising:} Suppose a single good is owned by a single seller who
values it at $c > 0$ (he can consume the good and get a payoff of $c$). There is a single buyer who
has a small transportation cost $k > 0$ to get to and back from the seller's store, and he values
the good at $v > c + k$. The buyer first decides whether to make the commute or stay at home, not
buy the good, and receive a payoff of 0. If the buyer commutes to the store, the seller can then make
the buyer a take-it-or-leave-it price offer $p \ge 0$. The buyer can then accept the offer, pay $p$,
and get the good, or he can walk out and not buy the good. Assume that $c$, $v$, and $k$ are common
knowledge.
\begin{partes}
\item As best you can, draw the extensive form of this game. What is the best response of the buyer
at the node at which he decides whether to accept or reject the seller's offer?
\item Find the subgame-perfect equilibrium of the game and show that it is unique. Is it Pareto
optimal?
\item Are there Nash equilibria that yield a higher payoff to both players as compared to the
subgame-perfect equilibrium you found in (b)?
\item Now assume that before the game is played the seller can, at a small cost
$\varepsilon < (v - c - k)$, send the buyer a postcard that commits the seller to a certain price
at which the buyer can buy the good (e.g., ``bring this coupon and get the good at a price $p$'').
Would the seller choose to do so? Justify your answer with an equilibrium analysis.
\end{partes}
\resolespaco

% ---------------------------------------------------------------------------
\exercicio[4]{Debt and Repayment}{Advanced}{Tadelis (2013), \emph{Game Theory: An Introduction}, Ch.~8
``Credibility and Sequential Rationality'', Exercise 8.7 ``Debt and Repayment'', p.~171 [PDF 190].
No answer in the book. Reproduced verbatim.}

\rot{8.7}\textbf{Debt and Repayment:} A project costing \$100 yields a gross return of \$110. A
lender (player 1) is approached by a debtor (player 2) requesting a standard loan contract to
complete the project. If the lender chooses \emph{not to offer} a loan, then both parties earn
nothing. If the lender chooses to \emph{offer} a loan of \$100, the debtor can realize the project's
gains and is obliged by contract to repay \$105. For simplicity assume that money is continuous and
that the debtor can choose to return any amount of money $x \le 110$. Ignore the time value of
money. Assume first that no legal system is in place that can compel the debtor to repay, so that
default on the loan (less than full repayment) carries no repercussions for the debtor.
\begin{partes}
\item Model this as an extensive-form game tree as best you can and find a subgame-perfect
equilibrium of this game. Is it unique?
\item Now assume that there is a legal system in place that allows the lender to
\emph{voluntarily} choose whether or not to sue when the debtor defaults and repays an amount
$x < 105$. Furthermore assume that it is \emph{costless} to use the legal system (it is supplied by
the state), and if the lender sues a debtor that defaulted, the lender will get the \$105 repaid in
full. After paying the lender, the borrower will pay a fine of \$5 to the court, above and beyond
the repayment. Model this as an extensive-form game tree as best you can and find a subgame-perfect
equilibrium of this game. Is it unique?
\item Are there Nash equilibria in the game described in (b) that are not subgame-perfect
equilibria?
\item Now assume that using the legal system is \emph{costly:} if the lender sues, he pays lawyers a
\emph{legal fee} of \$105 (this is the lawyers' price; it is unrelated to the contract). The rest of
the game proceeds as before. (If the lender sues a debtor that defaulted, the lender will get repaid
in full; after paying the lender, the borrower will pay a fine of \$5 above and beyond the
repayment.) Model this as an extensive-form game tree as best you can and find a subgame-perfect
equilibrium of this game. Is it unique?
\item Are there Nash equilibria in the game described in (d) that are not subgame-perfect
equilibria?
\item Now assume that a change in the law is proposed: upon default, if a debtor is sued he has to
first repay the lender \$105 and then pay the legal fees of \$105 above and beyond repayment of the
loan, and no extra fine is imposed. Should the lender be willing to pay to have this change in the
law enacted? If so, how much?
\item If you were a social planner, would you implement the suggested law?
\end{partes}
\resolespaco

% ---------------------------------------------------------------------------
\exercicio[5]{Observation, multiple SPNE and moving first}{Advanced}{Mas-Colell, Whinston \& Green
(1995), \emph{Microeconomic Theory}, Ch.~9 ``Dynamic Games'', Exercises, Exercise
9.B.5$^{\mathrm{B}}$ (E.~Maskin), p.~302 [PDF 162, left half]. Answer in the MWG solutions manual
(\texttt{Livros/}; 9.B.5 is there). Reproduced verbatim; the label of the last part is blurred in
the scan and is read as (e), as in the solutions manual.}

\rot{9.B.5$^{\mathrm{B}}$}(E. Maskin) Consider a game with two players, player 1 and player 2, in
which each player $i$ can choose an action from a finite set $M_i$ that contains $m_i$ actions.
Player $i$'s payoff if the action choices are $(m_1, m_2)$ is $\phi_i(m_1, m_2)$.

\textbf{(a)} Suppose, first, that the two players move simultaneously. How many strategies does each
player have?

\textbf{(b)} Now suppose that player 1 moves first and that player 2 observes player 1's move before
choosing her move. How many strategies does each player have?

\textbf{(c)} Suppose that the game in (b) has multiple SPNEs. Show that if this is the case, then
there exist two pairs of moves $(m_1, m_2)$ and $(m_1', m_2')$ (where either $m_1 \neq m_1'$ or
$m_2 \neq m_2'$) such that either
\[ \text{(i)}\quad \phi_1(m_1, m_2) = \phi_1(m_1', m_2') \]
or
\[ \text{(ii)}\quad \phi_2(m_1, m_2) = \phi_2(m_1', m_2'). \]

\textbf{(d)} Suppose that for any two pairs of moves $(m_1, m_2)$ and $(m_1', m_2')$ such that
$m_1 \neq m_1'$ or $m_2 \neq m_2'$, condition (ii) is violated (i.e., player 2 is never indifferent
between pairs of moves). Suppose also that there exists a pure strategy Nash equilibrium in the
game in \textbf{(a)} in which $\pi_1$ is player 1's payoff. Show that in any SPNE of the game in
\textbf{(b)}, player 1's payoff is at least $\pi_1$. Would this conclusion necessarily hold for any
\emph{Nash} equilibrium of the game in \textbf{(b)}?

\textbf{(e)} Show by example that the conclusion in \textbf{(d)} may fail either if condition (ii)
holds for some strategy pairs $(m_1, m_2)$, $(m_1', m_2')$ with $m_1 \neq m_1'$ or
$m_2 \neq m_2'$ or if we replace the phrase \emph{pure strategy Nash equilibrium} with the phrase
\emph{mixed strategy Nash equilibrium}.
\resolespaco
